\begin{afterword}
\summaryhead

The post starts from the problem of justifying Occam's razor. Arguing that the razor has worked before assumes the razor, since the hypothesis that it stops working after today fits the same record; calling a simplicity prior ``reasonable'' assumes it too. The author pictures a philosopher who may give up on justification and, by a truce with colleagues, label such beliefs ``a priori truth''.

Against this the post argues that thoughts are physical brain events. Adding 1 and 1 by ``pure thought'' is, in effect, observing your own brain, and watching another brain do the same computation would teach exactly the same thing. So nothing known ``a priori'' could not be known equally well by observation. The brain's bias toward simplicity is not arbitrary, as regularized learning algorithms show. Occam's razor works because the world is simple, though why it is simple remains unexplained. A mind that lacks Occam's razor or modus ponens cannot be argued into them, any more than a rock can. Brains were built by natural selection, not argued into existence, so rationality should be understood as engines that work, not arguments that convince.

\responsehead

The problem the post starts from is real and old. That the razor can be defended only by the razor is Hume's circularity argument about induction, applied to simplicity. The rival hypothesis that the razor fails after a given date has the form of Goodman's ``grue'', a predicate that changes its application at a set date. The complaint about a priori defences also has support: the Stanford Encyclopedia's entry on simplicity says that with them it can be hard to tell ``an a priori defense'' from ``no defense''. The portrait of philosophers calling a truce is offered with ``may'' and names no one. The same entry reports that many philosophers had turned away from such defences in the second half of the twentieth century.

The central argument has a gap. From the true premise that thoughts are brain events, the post concludes that knowing by ``pure thought'' is ``in effect'' observing one's brain. The definitions it quotes do not concern whether thinking is physical. Hume's independence from ``what is any where existent'' is independence from the objects the propositions are about, as his next sentence about circles and triangles shows, and the encyclopedia definition concerns what the belief's justification depends on. A person adding 1 and 1 has no observations of neurons among their reasons. The step works only if justification lies in the process that produces a belief. That is reliabilism, a position the post does not name, and reliabilists have used it to defend a priori knowledge. The equivalence with watching another brain can also run the other way: on one view of testimony, a belief received from someone who proved a theorem inherits the prover's a priori justification. So the conclusion that nothing is known a priori, which is one side of a live debate in philosophy, is not reached.

The hunter-gatherer example targets a claim the views it criticizes do not make. Their examples of a priori knowledge are arithmetic, logic and conceptual truths; nothing in them concerns where to find water.

The post's positive picture has good precedents. Judging reasoning by whether it works rather than by whether it can be argued for is the externalist reply to Hume's circularity: a method is justified if it is in fact reliable. The point about modus ponens is Lewis Carroll's, which the author later discussed. The claim that the brain is biased toward simplicity, given without a source, has support in cognitive science.

In short: a fair statement of an old problem, and an answer to it, judging methods by whether they work, with good precedents, together with an argument against the a priori that depends on a step from how a belief is produced to what justifies it, the same step reliabilists have used to defend the a priori.
\end{afterword}
